% Options for packages loaded elsewhere
\PassOptionsToPackage{unicode}{hyperref}
\PassOptionsToPackage{hyphens}{url}
\documentclass[
  ignorenonframetext,
]{beamer}
\newif\ifbibliography
\usepackage{pgfpages}
\setbeamertemplate{caption}[numbered]
\setbeamertemplate{caption label separator}{: }
\setbeamercolor{caption name}{fg=normal text.fg}
\beamertemplatenavigationsymbolsempty
% remove section numbering
\setbeamertemplate{part page}{
  \centering
  \begin{beamercolorbox}[sep=16pt,center]{part title}
    \usebeamerfont{part title}\insertpart\par
  \end{beamercolorbox}
}
\setbeamertemplate{section page}{
  \centering
  \begin{beamercolorbox}[sep=12pt,center]{section title}
    \usebeamerfont{section title}\insertsection\par
  \end{beamercolorbox}
}
\setbeamertemplate{subsection page}{
  \centering
  \begin{beamercolorbox}[sep=8pt,center]{subsection title}
    \usebeamerfont{subsection title}\insertsubsection\par
  \end{beamercolorbox}
}
% Prevent slide breaks in the middle of a paragraph
\widowpenalties 1 10000
\raggedbottom
\AtBeginPart{
  \frame{\partpage}
}
\AtBeginSection{
  \ifbibliography
  \else
    \frame{\sectionpage}
  \fi
}
\AtBeginSubsection{
  \frame{\subsectionpage}
}
\usepackage{iftex}
\ifPDFTeX
  \usepackage[T1]{fontenc}
  \usepackage[utf8]{inputenc}
  \usepackage{textcomp} % provide euro and other symbols
\else % if luatex or xetex
  \usepackage{unicode-math} % this also loads fontspec
  \defaultfontfeatures{Scale=MatchLowercase}
  \defaultfontfeatures[\rmfamily]{Ligatures=TeX,Scale=1}
\fi
\usepackage{lmodern}
\ifPDFTeX\else
  % xetex/luatex font selection
\fi
% Use upquote if available, for straight quotes in verbatim environments
\IfFileExists{upquote.sty}{\usepackage{upquote}}{}
\IfFileExists{microtype.sty}{% use microtype if available
  \usepackage[]{microtype}
  \UseMicrotypeSet[protrusion]{basicmath} % disable protrusion for tt fonts
}{}
\makeatletter
\@ifundefined{KOMAClassName}{% if non-KOMA class
  \IfFileExists{parskip.sty}{%
    \usepackage{parskip}
  }{% else
    \setlength{\parindent}{0pt}
    \setlength{\parskip}{6pt plus 2pt minus 1pt}}
}{% if KOMA class
  \KOMAoptions{parskip=half}}
\makeatother
\setlength{\emergencystretch}{3em} % prevent overfull lines
\providecommand{\tightlist}{%
  \setlength{\itemsep}{0pt}\setlength{\parskip}{0pt}}
% Slide layout defaults for warning-clean scientific decks.
\usepackage{etoolbox}
\IfFileExists{xurl.sty}{\usepackage{xurl}}{}
\IfFileExists{seqsplit.sty}{\usepackage{seqsplit}}{\newcommand{\seqsplit}[1]{#1}}
\protected\def\breakseq#1{\seqsplit{#1}}
\protected\def\breaktt#1{\begingroup\ttfamily\seqsplit{#1}\endgroup}
\setlength{\emergencystretch}{6em}
\tolerance=5000
\hbadness=10000
\hfuzz=1pt
\setlength{\tabcolsep}{2pt}
\AtBeginEnvironment{longtable}{\tiny\renewcommand{\arraystretch}{0.86}\setlength{\tabcolsep}{1pt}}
\AtBeginEnvironment{tabular}{\tiny\renewcommand{\arraystretch}{0.86}\setlength{\tabcolsep}{1pt}}

% Natbib and cross-reference fallbacks — slides don't load natbib
% or cleveref, but combined-PDF manuscript prose may emit these
% commands. The fallback renders citations as a bracketed key list
% and cross-references as detokenized labels so slides don't fail on
% undefined control sequences or raw underscores. \providecommand is
% a no-op if packages load later.
\providecommand{\citep}[1]{[#1]}
\providecommand{\citet}[1]{#1}
\providecommand{\citealp}[1]{#1}
\providecommand{\citeauthor}[1]{#1}
\providecommand{\citeyear}[1]{#1}
\providecommand{\cref}[1]{\texttt{\detokenize{#1}}}
\providecommand{\Cref}[1]{\texttt{\detokenize{#1}}}
\usepackage{bookmark}
\IfFileExists{xurl.sty}{\usepackage{xurl}}{} % add URL line breaks if available
\urlstyle{same}
% fallback for those not using the hyperref driver hyperxmp:
\makeatletter
\@ifundefined{xmpquote}{\newcommand{\xmpquote}[1]{#1}}{}
\makeatother
\hypersetup{
  hidelinks,
  pdfcreator={LaTeX via pandoc}}

\author{\texorpdfstring{}{}}
\date{}

\begin{document}

\begin{frame}[fragile,allowframebreaks]{Comparison to Existing Tools}
\protect\phantomsection\label{comparison-to-existing-tools}
The \href{02_introduction.md\#the-gap}{gap analysis} established that no
single tool integrates all six cross-cutting concerns. Here we
synthesize the
\href{08f_appendix_matrix.md\#appendix-matrix}{fourteen-dimension
comparison} into three structural insights. First, the landscape
bifurcates: workflow managers (Snakemake \breakseq{[@koster2012snakemake]},
Nextflow \breakseq{[@ditommaso2017nextflow]}, CWL \breakseq{[@amstutz2016cwl]})
provide distributed execution but no manuscript support; publication
tools (Quarto \breakseq{[@allaire2024quarto]}, Jupyter Book, R Markdown
\breakseq{[@xie2018dynamic]}, Overleaf \breakseq{[@overleaf2025]}, Prism
\breakseq{[@openai2026prism]}) author documents but embed no integrity
guarantees; and DVC \breakseq{[@iterative2024dvc]} versions artifacts without
orchestrating pipelines. \texttt{template/} occupies the intersection,
sacrificing distributed execution for unified enforcement of testing,
provenance, and documentation. This positioning is complementary---a
mature deployment might use Nextflow upstream and \texttt{template/} for
rendering, testing enforcement, and provenance downstream. Typst
\breakseq{[@madje2023typst]}, with its faster compilation cycle, is not one of
the nine compared peers but could serve as an alternative rendering
backend if a Pandoc writer were contributed.

Second, the eight enforcement capabilities \texttt{template/}
co-enforces---testing enforcement, coverage thresholds, steganographic
watermarking, multi-project management, AI-agent documentation, the
agentic skill protocol, an interactive TUI, and Zero-Mock policy---are
individually straightforward (and several, such as multi-project
management and AI-agent documentation, are matched in part by individual
peers); their novelty lies in \emph{co-enforcement} within a single
pipeline, ensuring that a passing build guarantees documentation
completeness, provenance embedding, and AI-navigability alongside
computational correctness. The FAIR4RS principles \breakseq{[@barker2022fair4rs;
@lamprecht2020towards]} articulate what research software quality
requires; FAIRsoft \breakseq{[@garijo2024fairsoft]} scores compliance
observationally; \texttt{template/} operationalizes these standards by
coupling them to pipeline gates that halt the build if coverage drops
below 90\% or provenance embedding fails. Cohen et al.'s four pillars of
research software engineering \breakseq{[@cohen2021four]}---sustainability,
quality, community, and policy---are operationalized by
\texttt{template/} through the first two pillars via quality-gated
automation.

Third, the AI-agent documentation dimension reveals an underserved need.
Overleaf and Prism provide AI \emph{writing} assistance, but neither
exposes structured documentation for \emph{external} agents to consume.
\texttt{template/}'s \texttt{AGENTS.md} + \texttt{SKILL.md} layer
enables an agent entering the repository to discover capabilities,
understand API contracts, and invoke modules without prior training
(\href{03c_documentation.md\#documentation-duality-and-ai-collaboration}{Documentation
Duality}, \href{05d_ai_collaboration.md\#the-ai-collaboration-model}{AI
Collaboration}).

\begin{block}{FAIR4RS Evolution (2024--2026)}
\protect\phantomsection\label{fair4rs-evolution-20242026}
Since the FAIR4RS principles were published \breakseq{[@barker2022fair4rs]},
the community has moved toward operationalization. The RDA Virtual
Plenary 24 (April 2025) featured a two-year retrospective review
\breakseq{[@honeyman2024fair4rs]} recommending principle amendments---notably
adding \emph{reproducibility} as an explicit requirement and clarifying
``domain-relevant standards''---alongside a leadership refresh and
parallel guidance activities. The ReSA Actionable FAIR4RS Task Force
(launched December 2024) analyzed the 17 principles into six actionable
categories (identifiers, metadata for publication/discovery/reuse,
standards, references, and licenses), with a first draft expected by
September 2025 \breakseq{[@resa2024actionable]}. Tools for automated FAIR
assessment have also matured: the F-UJI extension for research software
evaluation now scores against FRSM-04 through FRSM-17 metrics,
complementing Garijo et al.'s FAIRsoft evaluator
\breakseq{[@garijo2024fairsoft]}. \texttt{template/}'s pipeline-enforced
quality gates---coverage thresholds, documentation completeness checks,
and provenance embedding---anticipate this operationalization trend by
implementing FAIR4RS not as a post-hoc assessment but as an
architectural invariant.

In Gentleman and Temple Lang's terminology \breakseq{[@gentleman2007research]},
\texttt{template/} is a \emph{research compendium} scaled to the
repository level---bundling not just one study's code and data but N
studies, with shared infrastructure, automated testing, and embedded
provenance. Nüst et al.'s executable research compendium (ERC)
\breakseq{[@nust2017containerization]} extends this vision with containerized
reproduction environments; \texttt{template/} complements
containerization by adding the testing enforcement, multi-project
management, and provenance embedding layers that ERCs do not address.
\end{block}
\end{frame}

\end{document}
